\documentclass[11pt,a4paper,sans]{moderncv}
\moderncvstyle{classic}
\moderncvcolor{blue}
\usepackage[scale=0.75]{geometry}

\name{Behzad}{Haki}
\address{Barcelona, Spain / Vancouver, Canada}
\phone[mobile]{+34~684306775 / +1~(778)~386~1367}
\email{behzad.haki@upf.edu}
\homepage{www.behzadhaki.com}
% \social[linkedin]{behzad-haki-141bb18b}
\social[github]{behzadhaki}

\begin{document}
\makecvtitle

\section{Overview}

\cvitem{}{
% I am writing to express my interest in the Specialist Support Technician position with the Music Technology Group. With a background in Electrical and Computer Engineering and a Master's degree in Sound and Music Computing, and currently pursuing a PhD on generative music, I am confident in my ability to contribute to the projects involved in this position. I bring substantial experience in research, software development, web development, hardware development and proficiency in dev-ops tools. I have experience with Linux based clusters, and have hands-on experience with Machine Learning environments like TensorFlow and PyTorch. In addition to collecting and managing vast data collections, I developed many tools for curation, exploration and analysis of large datasets. Moreover, I have experience with deploying Neural Network models on embedded systems. My background in VST development with C++/JUCE will further complement the team's objectives. Lastly, I have extensive teaching and supervision experience. I believe that my ability to deliver autonomously, while synergizing as part of a team, makes me a suitable candidate for this role. 
}
% \cventry{}{}{}{}{}{}

% _____________________________________
\section{Education}

\cventry{2020--2024}{PhD in Sound and Music Computing}{\newline Universitat Pompeu Fabra (UPF)}{Barcelona, Spain}{}{
My dissertation, \textit{Design, Development, and Deployment of Real-Time Drum Accompaniment Systems}, examines the generation of real-time symbolic drum accompaniments with a focus on live improvisation contexts. Over the course of this research, I designed, developed, and evaluated three accompaniment systems of increasing complexity, each built around lightweight generative models optimized for real-time performance. A key aspect of this work was a close collaboration with professional musician Raül Refree, whose ongoing involvement shaped the design requirements and led to a series of live public performances with the developed systems. Beyond the primary systems, this work produced several secondary contributions, including NeuralMidiFx, a wrapper for deploying neural networks as VST plugins, two novel datasets (TapTamDrum and El Bongosero), and an exploration of audio-domain adaptations. A hardware Eurorack version of one of the systems was also designed and deployed.
}

\cventry{2017--2018}{Master in Sound and Music Computing}{\newline Universitat Pompeu Fabra (UPF)}{Barcelona, Spain}{}{
My master's studies at the Music Technology Group covered music perception and cognition,
real-time interaction, audio programming, digital signal processing, and deep learning.
Under the supervision of Sergi Jordà, I completed a thesis on generating basslines
interlocked with drum patterns using LSTM networks.
}

\cventry{2008--2013}{BASc in Electrical Engineering}{\newline University of British Columbia (UBC)}{Vancouver, Canada}{}{
I completed a degree in Applied Science with a major in Electrical and Computer
Engineering, specializing in Digital Signal Processing (DSP) and acoustics during
my final year.
}

% \cventry{2017--2018}{Master in Sound and Music Computing}{\newline Universitat Pompeu Fabra (UPF)}{Barcelona, Spain}{}{
% In the course of my master's studies, I specialized in sound and music computing at the Music Technology Group. Under the mentorship of Sergi Jordà, I immersed myself in a comprehensive curriculum that spanned music perception and cognition, real-time interaction, audio programming, digital signal processing, and deep learning. Each of these topics provided me with a holistic understanding of the intricate interplay between music and technology. For my thesis, I embarked on a project utilizing LSTM networks to generate basslines intricately interlocked with drum patterns.
% }
% % \cventry{}{}{}{}{}{}

% \cventry{2008--2013}{BASc in Electrical Engineering}{\newline University of British Columbia (UBC)}{Vancouver, Canada}{}{
% During my undergraduate years, I pursued a degree in Applied Science, majoring in Electrical and Computer Engineering. My interests began to converge on the fields of Digital Signal Processing (DSP) and acoustics. In my final year, I channeled my efforts and studies into these domains, delving deep into signal processing and the science of sound. This specialized focus laid the groundwork for my understanding of the intersection between engineering principles and auditory phenomena, equipping me with a robust foundation that would later influence my advanced academic pursuits.
% }
% % \cventry{}{}{}{}{}{}

% _____________________________________
\section{Skills}
\cvitem{Programming}{Software Design, Software Development, Software Testing, Python, C/C++, Matlab, Pure Data, Max/MSP, Concurrent Programming, Multi-Threaded Programming, PyTorch, Libtorch, ONNX, TensorFlow, JUCE, Docker, Linux, VST Plugin Development, Interface Development, Deployment of Neural Networks, Essentia.}
\cvitem{}{}

\cvitem{Hardware}{PCB Design, PCB Manufacturing, Circuit Design, Testing, Fabrication, Eurorack Module Development.}
\cvitem{}{}

\cvitem{Data Science}{Deep Learning. NLP Techniques for Music Generation. Collection, Curation and Processing of Large-scale Datasets. Data Analysis and Visualization. }
\cvitem{}{}

\cvitem{Acoustics}{EASE, CadnaA, ARTA, winMLS, LEAP EnclosureShop and CrossoverShop.}
\cvitem{}{}



% _____________________________________

\section{Work Experience}

\cventry{2024--2025}{Research Engineer}{Universitat Pompeu Fabra}{Barcelona, Spain}{}{
• Lead and conducted research on 
}
\cventry{}{}{}{}{}{}


\cventry{2014--2019}{Acoustic Engineer}{BAP Acoustics}{Vancouver, Canada}{}{
• Conduct noise and vibration measurements
• Model and simulate noise emissions from existing and proposed future noise sources
• Conduct room acoustic measurements
• Model and simulate acoustics of indoor spaces
• Simulate loudspeaker drivers based on electrical and acoustical measurements
• Design loudspeaker cabinets and cross-over circuits for specialized applications
• Simulate outdoor public warning systems
• Develop specialized software implementing acoustic standards
• Measure noise from railways to monitor the corrugation of tracks
• Write memorandums and reports
}
\cventry{}{}{}{}{}{}

% _____________________________________
\section{Teaching Experience}
\cventry{Fall 2019}{Introduction To Programming (First Year Engineering)}{\newline Department of Information and Communications Technologies (DTIC)
\newline Universitat Pompeu Fabra (UPF)}{}{}{This introductory course provides engineering students with a foundational understanding of programming, with a specific focus on Python. Through hands-on exercises, students are acquainted with the basics of Python programming, equipping them with essential coding skills that serve as a stepping stone for more advanced computational tasks.}
\cventry{}{}{}{}{}{}



\cventry{Winter 2020}{Electronic Music Production Lab (Fourth Year Engineering)}{\newline Department of Information and Communications Technologies (DTIC)
\newline Universitat Pompeu Fabra (UPF)}{}{}{Tailored for senior engineering students, this lab-oriented course delves into the realm of electronic music production. Utilizing the Pure Data environment, participants explore the intricacies of audio digital signal processing (DSP), gaining hands-on experience in crafting sound and understanding the underlying technical processes.}
\cventry{}{}{}{}{}{}

\cventry{Winter 2021, 2022, 2023}{Computer Organization  (First Year Engineering)}{\newline Department of Information and Communications Technologies (DTIC)}{}{}{A fundamental course for first-year engineering students, Computer Organization offers a deep dive into the architecture and operation of computers. Key topics include memory management techniques, the nuances of assembly language, and other foundational concepts that underpin the organization and functioning of computing systems.}
\cventry{}{}{}{}{}{}


\cventry{Winter 2021, 2022, 2023}{Computational Music Creativity (Masters)}{\newline Department of Information and Communications Technologies (DTIC)
\newline Universitat Pompeu Fabra (UPF)}{}{}{Situated at the intersection of music and technology, this master's level course delves into the world of computational music creativity. Participants are introduced to the basics of deep generative models, providing insights into how advanced algorithms can be leveraged to foster musical creativity and innovation.}
\cventry{}{}{}{}{}{}

% _____________________________________
% Publications
\nocite{*}
\bibliographystyle{unsrt}
\bibliography{publications}
\cventry{}{}{}{}{}{}




% _____________________________________
\section{Supervisions}
\cvitem{}{Throughout my academic journey at UPF, I have had the distinct privilege of supervising and mentoring a diverse group of students, aiding them in both their academic and research pursuits. My mentorship has spanned seven master's students from UPF, two undergraduates from the institution, as well as two visiting master's students and one visiting undergraduate. Notably, while it's often uncommon for master's students to achieve this feat, each of my master's mentees successfully published their works at esteemed conferences, a testament to our collaborative efforts and their dedication. These supervisory roles, both within and outside UPF, have been deeply enriching, underscoring my commitment to fostering academic excellence and pushing the boundaries of collaborative research.}


\cventry{}{}{}{}{}{}

\cventry{}{$\triangle$ Master's Students in Sound and Music Computing at MTG, UPF}{}{}{}{}
\cventry{}{}{}{}{}{}

\cventry{2022-23}{Nicholas Evans}{}{Deploying the GrooveTransformer to an Embedded Environment. \url{https://doi.org/10.5281/zenodo.8380628}}{}{}
\cventry{}{}{}{}{}{}

\cventry{}{Julian Lenz}{}{Disentangle and Deploy: Generative Rhythmic Tools for Musicians. \url{https://doi.org/10.5281/zenodo.8380515}}{}{}
\cventry{}{}{}{}{}{}

\cventry{2020-21}{Cheuk Lun Isaac Lee}{}{Dualisation of Rhythmic Patterns. \url{https://doi.org/10.5281/zenodo.5553940}}{}{}
\cventry{}{}{}{}{}{}

\cventry{}{Teresa Pelinski Ramos}{}{Completing Audio Drum Loops with Transformer Neural Networks. \url{https://doi.org/10.5281/zenodo.5554854}}{}{}
\cventry{}{}{}{}{}{}

\cventry{}{Marina S. Nieto Giménez}{}{Tap2Drum with Transformer Neural Networks. \url{https://doi.org/10.5281/zenodo.5554741}}{}{}
\cventry{}{}{}{}{}{}

\cventry{2019-20}{Błażej Kotowski}{}{Dualization Of Rhythm Patterns. \url{https://doi.org/10.5281/zenodo.4091469}}{}{}
\cventry{}{}{}{}{}{}

\cventry{}{Thomas Nuttall}{}{Transformer Neural Networks for Automated Rhythm Generation. \url{https://doi.org/10.5281/zenodo.4091477}}{}{}
\cventry{}{}{}{}{}{}
\cventry{}{}{}{}{}{}

\cventry{}{$\triangle$ Undergraduate Students}{}{}{}{}
\cventry{}{}{}{}{}{}

\cventry{2022-23}{Nil Distefano Capdeferro}{}{Drum Accompaniment for Instrumental Performances.}{}{}
\cventry{}{}{}{}{}{}

\cventry{}{Dídac Freire Castañé}{}{Drum Pattern Dualization.}{}{}
\cventry{}{}{}{}{}{}

\cventry{2019-20}{Aleix Rué Vilà}{}{Drum generation from a tapped pattern.}{}{}
\cventry{}{}{}{}{}{}

% \cventry{2022-23}{3 Master Students, 2 Undergraduate Student}{UPF, DTIC}{}{}{}
% \cventry{}{}{}{}{}{}

% \cventry{2020-21}{3 Master Students}{UPF, DTIC}{}{}{}
% \cventry{}{}{}{}{}{}

% \cventry{2019-20}{3 Master Students, 1 Undergraduate Student}{UPF, DTIC}{}{}{}






% _____________________________________
\section{Conference Reviewer}


\cventry{2025}{International Society for Music Information Retrieval Conference (ISMIR)}{}{}{}{}
\cventry{}{}{}{}{}{}

\cventry{2025}{International Computer Music Conference (ICMC)}{}{}{}{}
\cventry{}{}{}{}{}{}

\cventry{2025}{New Interfaces for Musical Expression Conference (NIME)}{}{}{}{}
\cventry{}{}{}{}{}{}

\cventry{2025}{International Conference on on AI and Musical Creativity (AIMC)}{}{}{}{}
\cventry{}{}{}{}{}{}

\cventry{2024}{New Interfaces for Musical Expression Conference (NIME)}{}{}{}{}
\cventry{}{}{}{}{}{}


\cventry{2023}{International Conference on on AI and Musical Creativity (AIMC)}{}{}{}{}
\cventry{}{}{}{}{}{}


\cventry{2020}{New Interfaces for Musical Expression Conference (NIME)}{}{}{}{}
\cventry{}{}{}{}{}{}




% _____________________________________
\section{Datasets}
\cventry{2023-24}{TapTamDrum}{}{\newline \url{https://taptamdrum.github.io/}}{}{A novel dataset of repeated dualizations from four experienced drummers. This collection consists of 1116 dualized drum patterns.}
\cventry{}{}{}{}{}{}

\cventry{2023-24}{El Bongosero}{}{\newline \url{https://elbongosero.github.io/}}{}{A large dataset of concurrent drum and bongo improvisations collected as part of the ongoing AI exhibition held at CCCB. This dataset consists of over 6000 improvisations from over 3000 participants.}
\cventry{}{}{}{}{}{}

% _____________________________________
\section{Open-Source Projects}
\cventry{2020-Now}{GrooveTansformer: A Suite of Generative Models}{}{\newline \url{https://github.com/behzadhaki/GrooveTransformer}}{}{This suite encompasses a collection of generative transformer models specifically tailored for the generation of symbolic drum patterns. Built on the transformer architecture, these models leverage advanced deep generative techniques to produce rhythmic sequences. The project is open-sourced and accessible to the global community, inviting collaboration and further exploration.}
\cventry{}{}{}{}{}{}



\cventry{2023}{GrooveTansformer: Eurorack Module}{}{}{}{In 2023, the 'GrooveTransformer' took a groundbreaking leap, evolving from a purely digital realm to the tactile world of Eurorack modules. Recognizing the potential of merging modern machine learning capabilities with the hands-on, modular approach of Eurorack, this initiative aimed to amplify user engagement and control in music generation. }
\cventry{}{}{}{}{}{}

\cventry{2023}{NeuralMidiFx}{}{\newline \url{https://neuralmidifx.github.io/}}{}{Addressing the deployment challenges of AI music generators in musician-favored platforms like DAWs, I developed NeuralMidiFx. This open-source template streamlines the integration of neural-based music systems as VST3 plugins. With a focus on ease of use, NeuralMidiFx reduces technical hurdles, empowering researchers to easily share their generative tools with a broader audience.}
\cventry{}{}{}{}{}{}

\cventry{2021-2023}{Groove2Drum VST}{}{\newline \url{https://github.com/behzadhaki/Groove2DrumVST}}{}{An advanced VST plugin deploying the GrooveTransformer modules, Groove2Drum offers musicians immediate generative drum pattern capabilities for both composition and accompaniment, seamlessly integrating AI-driven rhythms into standard music production workflows.}
\cventry{}{}{}{}{}{}

\cventry{2020-2023}{TapTamDrum Dataset}{}{\newline \url{https://taptamdrum.github.io/}}{}{Inspired by traditional drumming practices, TapTamDrum offers a dual-voice reduction dataset, capturing the essence of complex drum patterns using hands-only interpretations. Collected from four experienced drummers, this resource simplifies multi-voiced sequences for rhythm analysis and generation, providing a unique lens for rhythmic exploration.}
\cventry{}{}{}{}{}{}

\cventry{2020-22}{MonotonicGrooveTransformer}{}{\newline \url{https://github.com/behzadhaki/MonotonicGrooveTransformer}}{}{Designed for live musical accompaniment, this system utilizes a transformer encoder model trained to generate drum patterns from rhythmic inputs. Despite its training on short samples, strategic design enables it to accompany extended musical sequences seamlessly in real-time.}
\cventry{}{}{}{}{}{}

\cventry{2020-22}{TransformerGrooveInfilling}{}{\newline \url{https://transformergrooveinfilling.github.io/}}{}{Utilizing Transformer-based models, this tool augments audio drum loops with fitting symbolic events. By converging raw audio with symbolic transcriptions through a spectral approach, a real-time VST plugin is provided, seamlessly integrating its generative strengths into live production.}

\cventry{}{}{}{}{}{}


% _____________________________________

\section{Workshops}
\cventry{2023}{NeuralMidiFx Workshop at AIMC 2023}{University of Sussex, UK}{}{}{In this workshop, participants were introduced to NeuralMidiFx, a JUCE-based VST3 wrapper designed to simplify the integration of AI-driven symbolic music generative models into plugins. Tailored for those unfamiliar with plugin development, attendees were walked through the entire deployment process, culminating in the creation of a VST3 plugin using a pre-trained generative model.}
\cventry{}{}{}{}{}{}

% _____________________________________
\section{Artist Collaborations}
\cventry{2022-Now}{Collaboration with Raül Refree}{Barcelona, Spain}{}{}{Teaming up with Raül Refree, we've been a preparing a series of live concerts that integrate the GrooveTransformer. This collaboration is aimed at showcasing and exploring the fusion of traditional musical expertise with AI-driven generative tools}
\cventry{}{}{}{}{}{}


% _____________________________________
\section{Showcases}
\cventry{2023-Now}{CCCB: El Bongosero}{Barcelona, Spain}{\newline \url{https://elbongosero.github.io/}}{}{At the Centre de Cultura Contemporània de Barcelona (CCCB), we are currently running an interactive installation as part of a six-month exhibition. In this dynamic setting, museum visitors actively collaborate with a generative system. Beyond mere interaction, the installation is designed as a live platform for data collection. As attendees contribute their input, they play a vital role in the continuous enrichment and evolution of the model, ensuring its refinement throughout the exhibition's ongoing duration}
\cventry{}{}{}{}{}{}

\cventry{2023}{Sonar+D Music Festival}{Barcelona, Spain}{\newline \url{https://sonar.es/es/actividad/project-area-music-and-sound}}{}{During the Sonar+D festival, we showcased the GrooveTransformer Eurorack module to the general public.}
\cventry{}{}{}{}{}{}

\cventry{2023}{+RAIN Film Festival}{Barcelona, Spain}{\newline \url{https://www.upf.edu/es/web/rainfilmfest}}{}{In this concert, Nicholas Evans, my colleague and co-developer of the GrooveTransformer Eurorack Module, performed with the developed module}
\cventry{}{}{}{}{}{}

\cventry{Feb 2024}{CCCB Live Show}{Barcelona, Spain}{}{}{The first performance of Raül Refree with the GrooveTransformer system. In this live show, Refree did an improvised performance with the system.}
\cventry{}{}{}{}{}{}


\cventry{Nov 2024}{Palau Güell Live Show}{Barcelona, Spain}{}{}{The second performance of Raül Refree with the last transformer-based generative system. In this live show, Refree plays the organ at the venue, and the generative system generates accompaniments that will be played on a secondary organ. In this live show, we explore whether and how a generative system can be used in a context for which it was not designed.}
\cventry{}{}{}{}{}{}
% _____________________________________


%\cventry{2013}{Engineering Intern}{Serene Audio}{Vancouver, Canada}{}{
%• Worked on a study comparing Bluetooth encoding (APTX vs. SBC)    
%}





% _____________________________________
\section{Non-Academic Publications}
\cvitem{}{Haki, Behzad, De Santis, Eric. “Good Communication in Restaurants: Acoustic Capacity”. BAP
Acoustics Online Blog. May, 2015. \url{https://bapacoustics.com/good-communication-in-restaurants-acoustic-capacity/}}

\cvitem{}{Haki, Behzad. “How Good is Bluetooth at its best?”. Serene Audio Online Blog. November, 2015.
\url{http://www.sereneaudio.com/blog/how-good-is-bluetooth-audio-at-its-best}}

\cventry{}{}{}{}{}{}

% _____________________________________
\section{Languages}
\cvitemwithcomment{English}{Fluent}{}
\cvitemwithcomment{Farsi}{Fluent}{}
% \cvitemwithcomment{Spanish}{Beginner}{}

\cventry{}{}{}{}{}{}






\end{document}